\clearpage
\section{Demo de diagramas de flujo en LaTeX}

Sí: LaTeX puede dibujar correctamente el símbolo de \textbf{documento},
incluido su borde inferior ondulado. En este cuadernillo se usa TikZ, por lo
que el diagrama queda en formato vectorial y conserva su nitidez al imprimir.

\subsection{Símbolos solicitados}

\begin{figure}[H]
  \centering
  \begin{tikzpicture}[flujo,node distance=18mm]
    \node[documento] (doc) {Documento};
    \node[operacionManual,right=of doc] (manual) {Operación manual};
  \end{tikzpicture}
  \caption{Reproducción vectorial de los símbolos de las imágenes de referencia.}
\end{figure}

El primer símbolo representa convencionalmente un \textbf{documento o
reporte}; no representa el inicio del algoritmo. El segundo, con forma de
trapecio, representa una \textbf{operación manual}. Para inicio y fin se suele
usar el terminador ovalado. Si la convención indicada por el docente exige el
documento en otra posición, el estilo \texttt{documento} permite colocarlo sin
perder la forma ondulada.

\clearpage
\vspace*{0pt}
\subsection{Ejemplo completo}

\begin{figure}[H]
  \centering
  \begin{tikzpicture}[flujo,node distance=5mm and 16mm]
    \node[iniciofin] (inicio) {Inicio};
    \node[entradaSalida,below=of inicio] (leer) {Leer nombre};
    \node[decision,below=of leer] (valido) {¿Tiene contenido?};
    \node[proceso,below=of valido] (saludo) {Preparar saludo};
    \node[documento,below=of saludo] (reporte) {Generar documento};
    \node[iniciofin,below=of reporte] (fin) {Fin};
    \node[operacionManual,right=23mm of valido] (captura) {Capturar nombre\\manualmente};

    \draw (inicio) -- (leer);
    \draw (leer) -- (valido);
    \draw (valido) -- node[left,draw=none,fill=white] {Sí} (saludo);
    \draw (saludo) -- (reporte);
    \draw (reporte) -- (fin);
    \draw (valido) -- node[above,draw=none,fill=white] {No} (captura);
    \draw (captura) |- (leer);
  \end{tikzpicture}
  \caption{Flujo de ejemplo con terminador, entrada/salida, decisión, proceso,
  operación manual y documento.}
\end{figure}

\subsection{Código mínimo reutilizable}

El estilo ya está definido en \texttt{cuadernillo.tex}. Para insertar el
símbolo ondulado en cualquier actividad basta con escribir:

\begin{lstlisting}[language=TeX]
\begin{tikzpicture}[flujo]
  \node[documento] (resultado) {Reporte final};
\end{tikzpicture}
\end{lstlisting}
\thispagestyle{fancy}
